\documentclass[preview]{standalone}

\usepackage{amsmath}
\usepackage{amssymb}
\usepackage{stellar}
\usepackage{definitions}

\begin{document}

\id{psicologia-teorie-emozioni}
\genpage

\section{Teorie delle emozioni}

\begin{snippetdefinition}{teoria-periferica-emozioni-definition}{Teoria periferica delle emozioni}
    La \emph{teoria periferica delle emozioni}, o teoria di James-Lange,
    sostiene che l'emozione derivi dal feedback corporeo. La percezione di uno
    stimolo causa un'attivazione del sistema nervoso autonomo, e la percezione
    di tale attivazione conduce all'esperienza di una specifica emozione.
\end{snippetdefinition}

\begin{snippetexample}{james-lange-rabbia-example}{Rabbia e reazione corporea}
    Secondo la teoria periferica, non si urla perché si è arrabbiati; ci si
    sente arrabbiati perché si percepisce la reazione corporea collegata
    all'urlare e all'attivazione viscerale.
\end{snippetexample}

\begin{snippetdefinition}{teoria-centrale-emozioni-definition}{Teoria centrale delle emozioni}
    La \emph{teoria centrale delle emozioni}, o teoria di Cannon-Bard, sostiene
    che uno stimolo emotigeno produca due reazioni simultanee: l'attivazione
    fisiologica e l'esperienza emotiva soggettiva. Le due reazioni sono
    parallele e non dipendono causalmente l'una dall'altra.
\end{snippetdefinition}

\begin{snippetexample}{cannon-bard-rabbia-example}{Rabbia e simultaneità}
    Secondo la teoria centrale, si urla contro qualcuno perché si è arrabbiati:
    lo stimolo emotigeno produce in parallelo l'esperienza soggettiva della
    rabbia e l'attivazione fisiologica che prepara all'azione.
\end{snippetexample}

\begin{snippet}{teoria-periferica-centrale-confronto}
    La teoria periferica e la teoria centrale sono contrapposte ma parziali. La
    prima assegna un ruolo primario al feedback corporeo; la seconda attribuisce
    al cervello il coordinamento simultaneo di esperienza emotiva e attivazione
    fisiologica. Entrambe colgono aspetti reali ma non esauriscono la complessità
    dell'emozione.
\end{snippet}

\begin{snippetdefinition}{teoria-due-fattori-emozioni-definition}{Teoria emotiva a due fattori}
    La \emph{teoria emotiva a due fattori}, proposta da Schachter, sostiene che
    l'esperienza emotiva dipenda da due componenti necessarie: l'attivazione
    fisiologica, o \emph{arousal}, e l'attribuzione causale di tale attivazione
    nel contesto.
\end{snippetdefinition}

\begin{snippet}{sequenza-due-fattori}
    Nella teoria a due fattori, l'attivazione fisiologica è generica e
    indifferenziata. Il soggetto valuta il proprio stato corporeo, cerca
    etichette verbali adatte e attribuisce un significato alla propria
    attivazione in base al contesto in cui si trova.
\end{snippet}

\begin{snippetdefinition}{teoria-cognitivo-attivazionale-definition}{Teoria cognitivo-attivazionale}
    La \emph{teoria cognitivo-attivazionale delle emozioni} sostiene che
    l'esperienza emotiva dipenda dal modo in cui il soggetto valuta le
    interazioni in corso con l'ambiente, anche quando tale valutazione avviene
    senza pensiero consapevole.
\end{snippetdefinition}

\begin{snippet}{lazarus-attribuzione}
    Secondo Lazarus, l'attribuzione causale può avvenire senza ricerca esplicita
    nell'ambiente. Se l'esperienza passata lega certe emozioni a determinate
    situazioni, il soggetto può interpretare rapidamente la propria attivazione
    sulla base di valutazioni implicite.
\end{snippet}

\begin{snippetdefinition}{appraisal-definition}{Appraisal}
    L'\emph{appraisal} è il processo di valutazione cognitiva attraverso cui una
    situazione viene interpretata in funzione dei propri interessi, scopi e
    significati.
\end{snippetdefinition}

\begin{snippet}{teorie-appraisal}
    Le teorie dell'appraisal sostengono che le emozioni sorgano da un atto di
    conoscenza e valutazione della situazione. Le emozioni cambiano quando
    cambiano i significati o i valori attribuiti a una situazione; per questo
    due individui possono provare emozioni diverse davanti allo stesso evento.
\end{snippet}

\begin{snippet}{controlli-valutazione-stimolo}
    Secondo Scherer, l'appraisal procede attraverso controlli di valutazione
    dello stimolo:
    \begin{enumerate}
        \item valutazione della novità e della discrepanza rispetto alle
        aspettative;
        \item valutazione della qualità edonica, cioè della piacevolezza o
        spiacevolezza dello stimolo;
        \item valutazione della pertinenza rispetto a bisogni e scopi;
        \item valutazione delle capacità di coping;
        \item valutazione della compatibilità con norme sociali e immagine di sé.
    \end{enumerate}
\end{snippet}

\begin{snippetdefinition}{coping-emozioni-definition}{Coping}
    Il \emph{coping} è la valutazione delle proprie capacità di far fronte a uno
    stimolo emotigeno. Può riguardare il controllo dell'evento, la gestione delle
    reazioni emotive, la preparazione all'azione o la preparazione alla difesa.
\end{snippetdefinition}

\end{document}
